% Target: rmp-iv-ground-space-2d
% Formula: |psi(X)> = sum_{s} X( {beta_r} ) C( prod_r A^{s_r} ) |s>,
% Ink: Kernel plane frame; the sheet and its boundary operator are two nested
%   house enclosures over the box-site lattice, with X on the outer contour.
% Boundary: Resolved by the public environment; open indices remain explicit.
% Source: paper TN-Review-main.tex line 2006; author source ImagesReview Section
%   IV/ImagesReview Section IV.tex lines 55-70
% Capabilities: kernel, plane-frame, regions, typed-ports
\tenkzkernel
\tndeclare{atom}{tnsheetsite}
  {skin=box, ports={n:virtual, e:virtual, s:virtual, w:virtual}}
\[
  |\psi(X)\rangle \;=\;
  % The inner contour is the retained sheet; the outer contour, stepped one
  % clearance out, is the arbitrary boundary operator X acting on the sheet's
  % open indices.  Every site is the source's plain box tensor with one
  % undirected upward physical leg.
  \begin{tenkz}[lattice={4x4}, frame=plane, physical=up]
    \tnsheetsite{} & \tnsheetsite{} & \tnsheetsite{} & \tnsheetsite{} \\
    \tnsheetsite{} & \tnsheetsite{} & \tnsheetsite{} & \tnsheetsite{} \\
    \tnsheetsite{} & \tnsheetsite{} & \tnsheetsite{} & \tnsheetsite{} \\
    \tnsheetsite{} & \tnsheetsite{} & \tnsheetsite{} & \tnsheetsite{}
    \tnmark[form=enclosure]{(1,1) .. (4,4)}{}
    \tnmark[form=enclosure, inset=1, label pos=e]{(1,1) .. (4,4)}{$X$}
  \end{tenkz}
\]
